% page 540 du fac-simile (image p-539 du scan)
% bloc 154.9 x 246.3 mm ; marge gauche 8.0 mm
\pgc{-12.700mm}{0.000mm}{
\l{\cel{3}532}
\l{}
\l{\sou{sintona}. (\sou{elektro}.) Apta emisar ocili elektrala, di freque-}
\l{\cel{3}so intersama.}
\l{}
\l{sinui. (\sou{aludante voyo, rivero, fluvio, e c.)}\cel{3}Turni e jiri}
\l{\cel{3}ne-inter-egala. - EFIS.}
\l{}
\l{sinuso. I. (anat.)\cel{2}Kavajo kavernetatra di ula osti, di ula}
\l{\cel{3}vaskuli. - II. (\sou{geom.)}\cel{2}Penpendiklo trasita de un ek la}
\l{\cel{3}extremaji di arko til la diametro trasita ye la altra}
\l{\cel{3}extremajo. -}
\l{}
\l{\sou{sinusoido}. (geom.) Kurvo plana qua havas kom equaciono,}
\l{\cel{3}koordinati rektangula. : y = sin. x.}
\l{}
\l{\sou{sioro}. Titulo quan onu donas, pro politeso, ad omna adulti,}
\l{\cel{3}sive turnante su a li, sive parolante pri li.FIS.}
\l{}
\l{\sou{sireno}. I. (\sou{mitol.}) Ento mi-muliera, mi-fisha, qua atraktis}
\l{\cel{3}la mar-voyajanti til an la rifi, per la dolceso di sua}
\l{\cel{3}kanti. (Existis anke sireni ucela). - II. (\sou{metaf}.) Mu-}
\l{\cel{3}liero tre seduktiva. - III. (\sou{tekn}.) Instrumento en qua}
\l{\cel{3}sprico di vaporo, o di aero kompresita, produktas sono}
\l{\cel{3}stridanta, qua uzesas kom signalo sur la navi. - DEFIRS.}
\l{}
\l{\sou{sirfo}. (\sou{zool}.)Genero de insekti diptera, tipo di la familio}
\l{\cel{3}\textquotedbl{}sirfidei\textquotedbl{}, qui vivas omnaloke sur la Terglobo. - L.}
\l{\cel{3}\sou{syrphus}.}
\l{}
\l{\sou{siringo}. Mikra pumpilo kunportebla, uzata por klisteri od}
\l{\cel{3}altra injekti. - EFIS.}
\l{}
\l{\sou{siroko}. (\sou{navig}.) Nomo donita, sur Mediteraneo, a vento varm-}
\l{\cel{3}ega de sud-westo. - DEFIS.}
\l{}
\l{\sou{siropo}. I. (\sou{farmac.}) Liquido viskoza, ek sukro koquita,}
\l{\cel{3}unionita a ca o ta substanco di qua la rolo esas parfum-}
\l{\cel{3}izar lu per igar lu drinkajo agreabla. - II. Stando liqui-}
\l{\cel{3}da di la konfituri ante kande li divenis kelke solida per}
\l{\cel{3}koldesko. - III. Stando liquida di sukro qua esas suko}
\l{\cel{3}densa ante lua solidesko per koldesko. - DEFIR.}
\l{}
\l{\sou{sis}.\cel{2}La cifro \textquotedbl{}6\textquotedbl{}.\cel{2}EF.}
\l{}
\l{\sou{sisar}. (\sou{netrans.}) Produktar la son-uno\cel{1}\sou{sss}\cel{1}o\cel{1}\sou{fff}, t.e. bruiso}
\l{\cel{3}ne-muzikala qua esas, fonetike, konsonanto. (Kom exemplo:}
\l{\cel{3}serpento sisas ; lokomotivo siflas per sua siflilo, e}
\l{\cel{3}sisas pro emisar vaporo ek la kaldiero o la cilindri, tra}
\l{\cel{3}robineti, ma sen siflar). -}
\l{}
\l{\sou{sismo}. (\sou{geol.})\cel{2}Movo bruska di la Ter-kortico ; sukuso qua}
\l{\cel{3}shanceligas sulo sur plu o min granda areo, e qua esas}
\l{\cel{3}apta eventigar deformuri permananta en la regiono koncer-}
\l{\cel{3}nata. - DEFIRS.}
\l{}
\l{\sou{sismografo}. Cienco pri la sismi e pri la movi di sulo.}
}
